% Fuentes: EXP-010 (EXP-010_bs256_sqrt_60ep, 3 semillas), CHG-013; experiments.csv / history (lr, matrices de confusión).
\section{Resultado final}
\label{sec:final}
El modelo final (EXP-010) es \texttt{Rescaling} $\rightarrow$ \texttt{Flatten} $\rightarrow$ \res{final-hidden}
[\texttt{Dense} $\rightarrow$ LeakyReLU] $\rightarrow$ softmax sobre imágenes de
\res{final-res}~$\times$~\res{final-res} (\res{final-params} parámetros), entrenado con Adam ($\beta_1$ =
\res{final-beta1}, \emph{learning rate} \res{final-lr}) y batch de \res{final-batch}. El Cuadro~\ref{tab:final}
compara el modelo final con el baseline: el \emph{mean recall} se multiplica por \res{gain-total-ratio} y la
\emph{accuracy} sube \res{gain-total-acc}~pp (Figura~\ref{fig:progress}). El notebook de la entrega reproduce
exactamente el entrenamiento de la semilla~1.

\begin{table}[h]
  \centering\small
  \safetable{generated/tables/baseline_vs_final}
  \caption{Baseline frente al modelo final (validación; media $\pm$ desviación típica de \res{final-seeds} semillas).}
  \label{tab:final}
\end{table}

\begin{figure}[!htbp]
  \centering
  \safegraphics[width=0.85\linewidth]{figures/progress.png}{Progreso por fase}
  \caption{\figcaption{progress}}
  \label{fig:progress}
\end{figure}

\textbf{Punto de parada.} La Figura~\ref{fig:stopping} (semilla~1) muestra cómo decide la validación: hasta la época
\res{stop-last-joint-epoch} train y validación evolucionan en paralelo; al estancarse la validación,
\emph{ReduceLROnPlateau} divide la \emph{learning rate} por 10 y llega la mejor época (\res{stop-best-epoch}).
Después, la \emph{accuracy} de train sigue subiendo (\res{stop-train-acc-best} $\rightarrow$ \res{stop-train-acc-end})
mientras la de validación se estanca (\res{stop-val-acc-best} $\rightarrow$ \res{stop-val-acc-end}); las dos
bajadas siguientes de la \emph{learning rate} (épocas \res{stop-later-lr-drops}) no la mejoran, así que
\emph{EarlyStopping} (paciencia \res{final-es-patience}) para en la época \res{stop-epoch} y se conserva la
\res{stop-best-epoch}. La brecha entre train y validación pasa de \res{stop-gap-best} a \res{stop-gap-end}~pp entre la
mejor época y la parada: seguir entrenando solo habría servido para memorizar.

\begin{figure}[!htbp]
  \centering
  \safegraphics[width=0.6\linewidth]{figures/stopping.png}{Punto de parada}
  \caption{\figcaption{stopping}}
  \label{fig:stopping}
\end{figure}

Todas las clases mejoran respecto al baseline (Figura~\ref{fig:final}), aunque de forma muy desigual. La más difícil
sigue siendo \res{final-worst-class} (recall \res{final-worst-recall}), que se confunde con los otros vehículos:
de sus imágenes de validación, una proporción de \res{final-conf-truck-small-car} se clasifica como Small car,
\res{final-conf-truck-bus} como Bus y \res{final-conf-truck-dump-truck} como Dump truck (media de las matrices de
confusión de las \res{final-seeds} semillas). Entre train y validación queda una brecha (\emph{accuracy} \res{final-train-acc} frente a \res{final-val-acc}).

La mejor época y cada decisión se eligieron con el mismo conjunto de validación con el que se miden, así que todas
las cifras son algo optimistas~\cite{cawley}. Además, la desviación típica entre semillas solo recoge la variabilidad
de la inicialización y del orden de los datos, no la del \emph{split}, que es fijo; y el recall de las clases pequeñas
(p.\,ej.\ \res{class-min-name}, con pocas imágenes de validación) es especialmente ruidoso.

\begin{figure}[!htbp]
  \centering
  \safegraphics[width=0.85\linewidth]{figures/per_class_recall.png}{Recall por clase}
  \caption{\figcaption{per_class_recall}}
  \label{fig:final}
\end{figure}